\section{A uniform divisor moment}\label{sec:divisors}

We need a divisor estimate for intervals that may be much shorter than
their distance from zero. Only divisors up to a specified size are counted:
for \(X\ge1\) and \(n\in\N\), put
\[
 d_X(n)=\#\{a\in\N:a\le X,\ a\mid n\}.
\]
In the descent, this count bounds how many larger numerators can
lead to the same smaller residue. We need to control the combined
contribution of residues in a small exceptional set. H\"older's inequality
turns an \(r\)th-moment bound into a bound proportional to
\(\delta^{1-1/r}\) for a set occupying a proportion \(\delta\) of
an interval. Choosing \(r\) large limits the cumulative weakening of
this bound through the \(O(\log S)\) levels of the descent.
The moment exponent in the following lemma can be arbitrarily large,
provided it is fixed before the parameter \(S\) tends to infinity.
The proof follows the prime-factor splitting method of
Erd\H{o}s \cite[Section~3]{Erdos1952Divisors}; the uniform bound
for this truncated count is proved in full here.

\begin{lemma}\label{lem:divisor}
For every fixed pair of real numbers \(D,r\ge1\), there is
\(S_0(D,r)\) such that the following holds for \(S\ge S_0(D,r)\).
Suppose that \(X,Y\) are real numbers and \(N\) is a nonnegative integer
satisfying
\[
 \frac{S}{2\log S}\le \log X\le DS,\qquad
 X^{1/2}\le Y\le X,\qquad N\le e^{DS}.
\]
Then
\begin{equation}\label{eq:divisor-moment}
 \sum_{1\le h\le Y}d_X(N+h)^r\le Y\exp(S^{1/4}).
\end{equation}
\end{lemma}

\begin{proof}
Write \(v=\log X\), \(E=D+1\), and
\[
 B=1+\log(ES/v).
\]
Every integer \(n=N+h\) under consideration satisfies
\(1\le n\le2e^{DS}\le e^{ES}\) for sufficiently large \(S\).
Moreover,
\begin{equation}\label{eq:divisor-B}
 1<B\le1+\log(2E\log S)=O_D(\log\log S).
\end{equation}
All estimates below are uniform in \(X,Y,N\) in the stated ranges.

\smallskip
\noindent\emph{Two elementary estimates.}
First, consider divisors of \(n\) formed from primes at least \(z\ge2\).
List these prime factors with multiplicity, giving distinct labels to
repeated occurrences. There are at most
\(L=\lfloor ES/\log z\rfloor\) occurrences, and a divisor at most \(X\)
uses at most \(v/\log z\) of them. Counting subsets can only overcount
divisors. With \(q=v/(ES)\in(0,1)\), their number is therefore at most
\begin{equation}\label{eq:divisor-large-primes}
 \sum_{0\le j\le v/\log z}\binom Lj
 \le q^{-v/\log z}(1+q)^L
 \le \exp\left(\frac{Bv}{\log z}\right).
\end{equation}
Here the middle inequality follows by inserting the weights \(q^j\);
the last uses \(Lq\le v/\log z\).

Split the prime factors at \(S^{1/2}\).
For each smaller prime there are at most \(1+ES/\log2\) possible
exponents in a divisor of \(n\). Applying
\eqref{eq:divisor-large-primes} to the remaining primes gives
\[
 \log d_X(n)
 \le S^{1/2}\log(1+ES/\log2)+\frac{2Bv}{\log S}.
\]
Since \(v\ge S/(2\log S)\), there is a function
\(\varepsilon_D(S)\to0\), independent of \(X,Y,N,h\), such that
\begin{equation}\label{eq:divisor-pointwise}
 d_X(n)^r\le \exp\bigl(r\varepsilon_D(S)v\bigr),
 \qquad
 \varepsilon_D(S)\ll_D
 \frac{(\log S)^2}{S^{1/2}}+\frac{\log\log S}{\log S}.
\end{equation}

Second, write \(\tau(d)\) for the number of positive divisors of \(d\).
For fixed \(r\) and \(3/4\le\sigma\le1\),
\begin{equation}\label{eq:divisor-local-factor}
 \log\left(1+\sum_{j\ge1}(j+1)^r p^{-j\sigma}\right)
 \le C_r p^{-\sigma}
\end{equation}
for every prime \(p\).
Indeed, the power series
\(\sum_{j\ge1}(j+1)^r x^{j-1}\) is bounded on
\(0\le x\le2^{-3/4}<1\).
We will also use the elementary estimate
\begin{equation}\label{eq:prime-reciprocals}
 \sum_{p\le T}\frac1p\le e\log(1+\log T)\qquad(T\ge2).
\end{equation}
To prove it, set \(s=1+1/\log T\) and
\(\zeta(s)=\sum_{a\ge1}a^{-s}\).
Euler's absolutely convergent product and the integral bound
\(\zeta(s)\le1+1/(s-1)\) give
\[
 \sum_{p\le T}\frac1p
 \le e\sum_p p^{-s}
 \le e\log\zeta(s)
 \le e\log(1+\log T).
\]
In particular, multiplicativity and
\eqref{eq:divisor-local-factor} imply
\begin{equation}\label{eq:divisor-harmonic}
 \sum_{d\le T}\frac{\tau(d)^r}{d}
 \le \prod_{p\le T}
       \left(1+\sum_{j\ge1}(j+1)^r p^{-j}\right)
 \ll_r(\log(2T))^{C_r}.
\end{equation}

\smallskip
\noindent\emph{A prefix of the prime factorization.}
For \(n>\sqrt Y\), order its prime factors with multiplicity and let
\(d\) be the longest initial product at most \(\sqrt Y\), allowing
\(d=1\). Let \(p\) be the next prime. Then
\[
 d\le\sqrt Y<dp,
\]
all prime factors of \(d\) are at most \(p\), and all those of \(n/d\)
are at least \(p\).
The inequality
\[
 d_X(n)\le \tau(d)d_X(n/d)
\]
does not require \(d\) and \(n/d\) to be coprime:
a divisor \(a\le X\) of \(n\) determines
\[
 b=\gcd(a,d),\qquad c=a/b.
\]
Then \(b\mid d\), \(c\mid n/d\), \(c\le X\), and the pair \((b,c)\)
determines \(a\).
Thus \eqref{eq:divisor-large-primes} gives
\begin{equation}\label{eq:divisor-prefix}
 d_X(n)^r\le\tau(d)^r
          \exp\left(\frac{rBv}{\log p}\right).
\end{equation}
For any fixed \(d\le\sqrt Y\), the number of its multiples in
\((N,N+Y]\) is at most
\begin{equation}\label{eq:divisor-multiples}
 \frac Yd+1\le\frac{2Y}{d}.
\end{equation}
We now sum according to the size of the next prime \(p\). If \(p\)
is large, the remaining factor \(n/d\) has few small divisors. If
\(p\) is smaller, the condition \(dp>\sqrt Y\) forces \(d\) to
be a large product of small primes. The total reciprocal weight of
such prefixes is small, so \eqref{eq:divisor-multiples} limits their
occurrence in every shifted interval.

If \(\log p\ge v^{15/16}\), the exponential factor in
\eqref{eq:divisor-prefix} is at most \(e^{rBv^{1/16}}\).
For \(n\le\sqrt Y\), instead take \(d=n\) and use
\(d_X(n)\le\tau(d)\). By \eqref{eq:divisor-multiples} and
\eqref{eq:divisor-harmonic}, these two cases together contribute at most
\begin{align}
 2Ye^{rBv^{1/16}}\sum_{d\le\sqrt Y}\frac{\tau(d)^r}{d}
 &\le Y\exp\!\left(
      O_{D,r}(S^{1/16}\log\log S+\log S)\right).
 \label{eq:divisor-large-contribution}
\end{align}

For the remaining integers, \(\log p<v^{15/16}\), so
\begin{equation}\label{eq:divisor-prefix-lower}
 \log d>\tfrac12\log Y-\log p
 \ge v/4-v^{15/16}>v/5
\end{equation}
once \(S\) is sufficiently large.
This large prefix makes integers with small next prime rare.
We estimate the ensuing smooth-prefix sums by Rankin's exponential
weighting method; see \cite[p.~414, (1.3)]{HildebrandTenenbaum1993}.

If \(p<S^4\), every prime factor of \(d\) is smaller than \(S^4\).
Multiplying each summand by
\((d/e^{v/5})^{1/10}>1\) gives
\[
 \sum_{\substack{d>e^{v/5}\\
             \ell\mid d,\ \ell\ {\rm prime}\Rightarrow\ell<S^4}}
       \frac1d
 \le e^{-v/50}
       \prod_{\substack{\ell<S^4\\\ell\ {\rm prime}}}
                      (1-\ell^{-9/10})^{-1}.
\]
The logarithm of this product is \(O(S^{2/5})\), since
\[
 \sum_{\substack{\ell<S^4\\\ell\ {\rm prime}}}\ell^{-9/10}
 \le\sum_{2\le a<S^4}a^{-9/10}=O(S^{2/5}).
\]
Counting multiples by \eqref{eq:divisor-multiples} and bounding the
entire summand by \eqref{eq:divisor-pointwise}, we obtain a contribution
at most
\begin{equation}\label{eq:divisor-small-contribution}
 2Y\exp\left(-v/50+O(S^{2/5})+
                         r\varepsilon_D(S)v\right)
 \le 2Ye^{-v/100}.
\end{equation}
Here \(S^{2/5}=o(v)\), uniformly in the allowed range.

\smallskip
\noindent\emph{The intermediate primes.}
It remains to treat
\(4\log S\le\log p<v^{15/16}\).
Partition this range into intervals \([t,2t)\), where
\(t=2^j4\log S\le v^{15/16}\).
There are \(O_D(\log S)\) such intervals; the last may extend beyond
the upper endpoint. For one interval put
\[
 Q=\frac vt,\qquad \lambda=\frac{\log Q}{10t}.
\]
Then \(Q\ge v^{1/16}\), and
\[
 0<\lambda
 \le\frac{\log(DS/(4\log S))}{40\log S}<\frac14
\]
for sufficiently large \(S\).
Equations \eqref{eq:divisor-prefix},
\eqref{eq:divisor-multiples}, and \eqref{eq:divisor-prefix-lower}
bound the contribution of this interval by
\[
 2Ye^{rBQ}
 \sum_{\substack{d>e^{v/5}\\
           \ell\mid d,\ \ell\ {\rm prime}\Rightarrow\ell\le e^{2t}}}
       \frac{\tau(d)^r}{d}.
\]
The last sum is at most
\[
 e^{-\lambda v/5}
 \prod_{\substack{\ell\le e^{2t}\\\ell\ {\rm prime}}}
 \left(1+\sum_{j\ge1}(j+1)^r\ell^{-j(1-\lambda)}\right).
\]
By \eqref{eq:divisor-local-factor} and
\eqref{eq:prime-reciprocals}, the logarithm of this product is
bounded by
\[
 C_r\sum_{\substack{\ell\le e^{2t}\\\ell\ {\rm prime}}}
                   \ell^{-1+\lambda}
 \ll_r e^{2t\lambda}\log(2t)
 =Q^{1/5}O_r(\log(2t)).
\]
The contribution of the interval is consequently at most
\begin{equation}\label{eq:divisor-band}
 2Y\exp\left(rBQ-\frac{Q\log Q}{50}
                         +O_r(Q^{1/5}\log(2t))\right).
\end{equation}
Both positive terms in this exponent are uniformly negligible compared
with \(Q\log Q\). Indeed, \eqref{eq:divisor-B} and
\(Q\ge v^{1/16}\) give
\[
 \frac{rB}{\log Q}
 \ll_{D,r}\frac{\log\log S}{\log S}\longrightarrow0,
 \qquad
 \frac{Q^{1/5}\log(2t)}{Q\log Q}
 \ll v^{-1/20}\longrightarrow0.
\]
Thus \eqref{eq:divisor-band} is at most
\(2Y\exp(-Q\log Q/100)\le2Y\), after increasing \(S_0(D,r)\).
Summing the intervals costs \(O_D(Y\log S)\).

Combining this bound with \eqref{eq:divisor-large-contribution}
and \eqref{eq:divisor-small-contribution} gives
\[
 \sum_{1\le h\le Y}d_X(N+h)^r
 \le Y\exp\!\left(
         O_{D,r}(S^{1/16}\log\log S+\log S)\right)
       +2Ye^{-v/100}+O_D(Y\log S).
\]
The logarithmic loss is \(o(S^{1/4})\), proving
\eqref{eq:divisor-moment}. Every enlargement of the threshold depended
only on the fixed \(D,r\); this completes the required uniformity check.
\end{proof}
